\section*{6.2 Gadget classification for randomized communication complexity (author source pp.39--40)}
In this section, we prove the following dichotomy result for the randomized case.
\begin{originalresult}{Theorem 6.7}
For every gadget $g$, one of the following cases holds:
\begin{enumerate}
\item[(i)] There is an infinite family of functions $f_n$ with $\operatorname{bs}(f_n)=\Omega(n)$ and $R(f_n\comp g)=O(\log n)$.
\item[(ii)] For every function $f$, we have $R(f\comp g)=\Omega(\operatorname{bs}(f)^{\Omega(1)})$.
\end{enumerate}
\end{originalresult}
First, we need to prove the following lemma.
\begin{originalresult}{Lemma 6.8}
Suppose that gadgets $g,h:X\times Y\to\bits$ are such that $\OR\nleq g$ and $\AND\nleq h$. Then
\[R(\AND_n\comp g)=O(\log n),\qquad R(\OR_n\comp h)=O(\log n).\]
\end{originalresult}
\begin{proof}
The bounds for $g$ and $h$ are symmetrical, so we prove only the first one. Consider the decomposition $X=X_0\sqcup X_1\sqcup\cdots\sqcup X_k$ and $Y=Y_0\sqcup Y_1\sqcup\cdots\sqcup Y_k$ provided by Lemma 6.3. If one of the blocks $g_i$ was given an input from $X_0$ or $Y_0$, then one of the players can share this information with the other player and finish the communication. Otherwise, they indicate that they do not have such blocks. After this step, they are restricted to inputs from $X\setminus X_0$ and $Y\setminus Y_0$. Since any two inputs from the same block $X_i$ or $Y_i$ are indistinguishable for our gadget, they need to solve $\AND_n\comp\operatorname{EQ}_m$, where $\operatorname{EQ}_m$ is the equality function for $m$-bit inputs (i.e., $\operatorname{EQ}_m(x,y)=1\iff x=y$) for $m=\lceil\log k\rceil$. This is exactly an instance of $\operatorname{EQ}_{nm}$, which can be solved with $O(\log n)$ bits of communication (note that $m$ is a constant) by a randomized communication protocol (see Kushilevitz \& Nisan (1996)).
\end{proof}
Now we are ready to prove the dichotomy result.
\begin{proof}[Theorem 6.7]
If $\AND,\OR\le g$ then Theorem 6.1 shows that $R(f\comp g)=\Omega(\operatorname{bs}(f)^{\Omega(1)})$, so we get the second case.
Otherwise, $\OR\nleq g$ or $\AND\nleq g$. We apply Lemma 6.8: if $\OR\nleq g$ we take $f_n=\AND_n$, otherwise we take $f_n=\OR_n$. This gives us the first case of the theorem.
\end{proof}
